\section{Next targets}\label{sec:21.7}

CsI has a heavy-mediator single-phonon rate per unit mass 3.73 times that of silicon, on elastic constants measured at \SI{4.2}{K}, and 1.56 times that of TeO$_2$, the leader of the eleven crystals ranked in \cref{sec:17.6}. The lead salts have electron-recoil floors up to 4.4 times below germanium's. Neither family was in that pool: no CsI crystal has been read by a phonon sensor at millikelvin temperatures, and a lead compound carries the $^{210}$Pb of its lead, whose $\beta$ spectrum puts some 700 events per kilogram, day and keV below \SI{1}{keV} for every becquerel per kilogram. A ranking of responses answers which crystal gives the most rate per unit mass. An experimentalist choosing the target of the next experiment also needs to know what each leader still has to demonstrate, which measurement would change the order, and on what the probabilities rest. We combine the response-level rankings of \cref{sec:17.6} with the pair energies of \cref{sec:18.4}, the readout verdicts of \cref{sec:20.5} and the internal backgrounds of the crystals into a guide with one row per channel, \cref{tab:21-guide}.

\subsection{Candidate pools and detector readiness}\label{sec:21.7-pools}

The \emph{full pool} holds 42 crystals: 23 of the 24 crystals of the target table of Griffin and collaborators~\cite[Table~I]{griffin2020multichannel}, all but InSb, the three SiC polytypes, and the targets of the preceding chapters that have the inputs of at least one figure of merit. Four requirements decide whether a crystal can be the target of a cryogenic experiment now. The crystal operates at millikelvin temperatures after thermal cycling, without a phase transition between \SI{300}{K} and \SI{10}{mK} and without absorbing water from the air. At least \SI{1}{g} of detector-grade crystal can be grown, or an array of instrumented crystals of at least \SI{0.1}{g} each adds up to a gram. A sensor has been read on its surface at millikelvin temperatures. And its intrinsic rate below \SI{1}{keV} is at most one event per kilogram, day and keV. A crystal for which all four are documented is a \emph{survivor}. A crystal for which one or more is not documented, though none is shown to be impossible, is a \emph{conditional survivor}. Only positive evidence removes a crystal from the full pool. That evidence exists for two candidates, KTaO$_3$ and InSb, whose natural $^{40}$K and $^{115}$In give intrinsic rates far above the fourth requirement (\cref{sec:21.7-backgrounds}).

The survivor pool is the eleven crystals of \cref{sec:17.6}: silicon, germanium, GaAs, diamond, sapphire, LiF, NaF, CaWO$_4$, TeO$_2$, Li$_2$MoO$_4$ and $\gamma$-LiAlO$_2$. The other 31 are conditional. Membership records what has been documented for a crystal, not its physics, and it changes with one demonstration. GaAs entered when a \SI{4.3}{g} crystal was operated at about \SI{10}{mK} with a thermistor readout~\cite{melchiorre2025first}. Diamond entered through the array route: two crystals of \SI{0.175}{g} each, grown by chemical vapor deposition and read by tungsten transition-edge sensors, one with a threshold of \SI{16.8}{eV}~\cite{angloher2023light}. Six such crystals make a gram, each with its own readout channel and threshold, and their surface-to-volume ratio is $(1/0.175)^{1/3} = 1.79$ times that of a single gram crystal, which raises any surface background in proportion. The survivor leader of a channel is the best target that the record supports today. The full-pool leader says what a demonstration would gain: if it differs from the survivor leader, the missing demonstration of the full-pool leader is worth as much as its lead.

The rankings of the full pool use the method and widths of \cref{sec:17.6} unchanged. For the conditional crystals the gaps are those of the target table, a room-temperature gap shifted upward by an amount drawn uniformly between 0.03 and \SI{0.10}{eV}, and the SiC gaps are room-temperature values of the SiC target table~\cite[Table~I]{griffin2021sic}. Their elastic tensors are computed ones from the same database as in \cref{sec:17.6}~\cite{choudhary2018elastic}, with the width 0.086, except two that are measured: aluminum at zero temperature and CsI at \SI{4.2}{K}~\cite{vallin1964elastic}, width 0.02. Their polar band integrals come from \cref{eq:11.2-diatomic} with computed Born charges, $\varepsilon_\infty$ and highest optical energies~\cite[Table~I]{griffin2020multichannel}, or from the measured dielectric constants of the SiC polytypes~\cite[Table~I]{griffin2021sic}.

We add one channel that \cref{sec:17.6} did not rank: the threshold in quanta, measured by the pair energy $\varepsilon_{\mathrm{eh}}$. A lower pair energy turns the same deposit into more pairs and puts the one- and two-pair bins of a charge readout at lower energy (\cref{sec:18.3}). Where a pair energy has been measured, the guide uses the measurement: silicon \SI{3.749}{eV} at \SI{123}{K}~\cite{rodrigues2021absolute}, GaAs about \SI{4.2}{eV} at room temperature~\cite[Fig.~1]{nielsen2026first} raised by $1.70 \times \SI{0.10}{eV}$ for the shift of its gap on cooling (the slope $\partial\varepsilon_{\mathrm{eh}}/\partial\Eg$ of the cascade at fixed constant, \cref{sec:18.4}), diamond and the SiC polytypes as in \cref{fig:18-limit-c}, and germanium as discussed below. For every other crystal the pair energy is the estimate of the cascade of \cref{sec:18.4} at the tabulated gap and highest optical phonon energy. Inside the validated range of \cref{sec:18.4-range} we give it a logarithmic width 0.105, the root mean square of the residuals of the seven crystals with measured values; outside it the width is doubled to 0.21, a convention, since no measurement anchors these crystals.

Germanium has no pair energy measured directly at millikelvin temperatures. Its pair energy at \SI{77}{K} is \SI{2.96}{eV}, the value \cref{sec:18.4} tests the cascade against. The one value quoted at millikelvin temperatures, $(3.37 \pm 0.01)$~eV, comes from a fit by Wei, Wang and Mei to the ratio of phonon to ionization energy measured in EDELWEISS detectors as a function of bias, $E_{\mathrm{p}}/E_{\mathrm{Q}} = a + eV_{\mathrm{bias}}/\varepsilon_{\mathrm{eh}}$, in which the relative conversion efficiency $a = 0.85$ of primary and Luke phonons is fitted together with $\varepsilon_{\mathrm{eh}}$~\cite[Sect.~2.1.4]{wei2016average}. The value rests on the assumption that one constant $a$ describes the phonon collection at every bias, the quantity whose detector dependence \cref{sec:20.5} found open. We use 3.37~eV with the width 0.015 of the spread between this fit and their resolution-based value of 3.32~eV. Either value puts germanium ahead of silicon.

Every probability in the guide is conditional on the provenance widths of \cref{sec:17.6-widths}. The guide reads every row at independent method errors, $r_{\mathrm{mat}} = 0$ in \cref{eq:17.6-zmethod}, and prints the value at $r_{\mathrm{mat}} = 1$ beside it.

\subsection{The guide to next targets}\label{sec:21.7-guide}

\Cref{tab:21-guide} is the guide. The survivor leaders and their values, except in the pair-energy row, are those of \cref{tab:17-rank}.

\begin{table}[tbp]
  \caption{The guide to next targets. Leaders of the full pool (42 crystals) and of the survivor pool (eleven crystals, \cref{sec:21.7-pools}). Full-pool leader: median with 68\% interval, and probability of ranking first at independent method errors with the value at $r_{\mathrm{mat}} = 1$ in parentheses; the values of the survivor leaders are in \cref{tab:17-rank}. Every probability is conditional on the provenance widths of \cref{tab:17-widths} and on the widths of the pair energies stated in \cref{sec:21.7-pools}; it is not the probability that a crystal is the better target. Missing demonstration: the requirements of \cref{sec:21.7-pools} not documented for the leader. Edges at $\omega_{\mathrm{th}} = \SI{1}{meV}$; $v_{\mathrm{max}} = 2.650\times10^{-3}$.}
  \label{tab:21-guide}
  \centering\scriptsize\setlength{\tabcolsep}{3pt}
  \begin{tabular}{@{}>{\raggedright\arraybackslash}p{0.14\linewidth}>{\raggedright\arraybackslash}p{0.24\linewidth}>{\raggedright\arraybackslash}p{0.12\linewidth}>{\raggedright\arraybackslash}p{0.22\linewidth}>{\raggedright\arraybackslash}p{0.20\linewidth}@{}}
    \toprule
    Channel; figure of merit & Full-pool leader & Survivor leader & Missing demonstration & Deciding measurement \\
    \midrule
    Electron-recoil floor; $2\Eg/v_{\mathrm{max}}^2$ & PbSe, \SI{0.0484}{MeV} [0.0479, 0.0489]; 1.000 (1.000); PbTe and PbS next & Ge & lead salts: a crystal of at least \SI{0.1}{g} read at mK; $^{210}$Pb below \SI{1.38}{mBq} per kg of compound & $^{210}$Pb activity of the lead available for growth; low-temperature gap of GaSb \\
    Threshold in quanta; $\varepsilon_{\mathrm{eh}}$ & PbTe, \SI{1.52}{eV} [1.24, 1.87]; 0.482 (1.000); one tier with PbSe (0.414) and PbS (0.102) & Ge, \SI{3.37}{eV} [3.32, 3.42]; then Si 3.75, GaAs 4.37 & lead salts: as above; one measured pair energy of a lead salt & pair energy of germanium measured at mK with the phonon collection calibrated separately \\
    Heavy-mediator single phonon; $c_{\mathrm{NR1}}(\mathrm{Si})/c_{\mathrm{NR1}}$ & CsI, 3.73 [3.65, 3.80]; 0.962 (0.976); measured \SI{4.2}{K} tensor & TeO$_2$, computed tensor & CsI: an encapsulated crystal cycled to mK and read by a phonon sensor & measured elastic constants of TeO$_2$; measured $^{123}$Te capture rate \\
    Nuclear edge; $m_{\mathrm{min}}^{\mathrm{ac}}$ with $c_{\mathrm{LA,max}}$ & no resolved leader: h-BN, \SI{2.95}{keV} [2.72, 3.21]; 0.656 (0.728); one tier with diamond (0.181), polyethylene (0.119), cubic BN (0.043) & diamond, measured tensor & diamond: six crystals of \SI{0.175}{g}, each with its own threshold and low-energy-excess rate, and a phonon threshold near \SI{1}{meV} (shown: \SI{16.8}{eV}); h-BN: mass, mK and readout & measured elastic constants of h-BN, polyethylene and cubic BN \\
    Polar rate; $\mathcal{I}_{\mathrm{pol}}/\rho_T$ & LiF, 19.0 [15.5, 23.3]; 0.918 (0.931) & LiF & LiF: the $^6$Li neutron-activation background, not measured & measured $\varepsilon_0$ and Born charge of LiF at low temperature \\
    Polar reach; $2\omega_{\mathrm{L}}/v_{\mathrm{max}}^2$ & CsI, \SI{2.85}{keV} [2.71, 2.99]; 0.996 (1.000); computed $\omega_{\mathrm{L}}$ & GaAs, measured $\omega_{\mathrm{L}}$ & GaAs: a threshold below its longitudinal phonon, \SI{36.35}{meV}; CsI: as above & low-temperature longitudinal optical frequency of CsI \\
    Migdal; pole factor times $X_{\mathrm{pole}}$ & no resolved leader: NaCl 0.498 (1.000) & no resolved leader: NaF 0.315 (1.000) & -- & $L(\vect{k},\omega)$ of every candidate \\
    \bottomrule
  \end{tabular}
\end{table}

Germanium leads both electron-recoil rows among the survivors, the floor on its measured low-temperature gap and the threshold in quanta with 1.11 times fewer electron volts per pair than silicon. Its lead in quanta holds on the \SI{77}{K} value as well, which is lower still, so the 50~mK fit affects the number and leaves the order unchanged. In the full pool the lead salts lead both rows. PbSe has the lowest floor, on a gap measured at low temperature (\cref{fig:17-rank-floor}). In quanta PbTe, PbSe and PbS share the leading tier. Their pair energies are cascade estimates for crystals whose phonon strengths $\lambda_{\mathrm{AB}}$, the dimensionless ratio of phonon emission to impact ionization in the cascade (\cref{sec:18.4-range}), lie between 230 and 1160, two orders of magnitude outside its validated range, so the tier is real only under the width convention. GaSb, whose estimate lies at the edge of that range, beats germanium's measured value with probability 0.83. At $r_{\mathrm{mat}} = 1$ the three estimates of the lead salts share their error and PbTe ranks first in every draw. No lead salt has been operated at millikelvin with a phonon or charge readout, and every one carries $^{210}$Pb in its bulk (\cref{sec:21.7-backgrounds}).

The heavy-mediator rows separate the two pools most clearly. TeO$_2$ leads the survivors on a computed tensor. CsI, a conditional survivor, leads the full pool on elastic constants measured at \SI{4.2}{K}, from which the Christoffel problem of \cref{sec:17.6} gives $c_{\mathrm{NR1}} = \SI{2.404}{km.s^{-1}}$ and the figure of merit $8.956/2.404 = 3.73$ times silicon. It is ahead of Bi$_2$Se$_3$ and NaI, both computed, with probabilities 0.978 and 0.985, and of TeO$_2$ by a factor 1.56 in the medians, larger than any width. A colder, denser cell, at most \SI{4.5}{\percent} denser than the room-temperature one, raises the figure of merit to 3.81. Against silicon the ratio of the figures of merit gives the ratio of the rates to \SI{10}{\percent} from \SI{97}{keV} up to the zone cap of CsI at \SI{0.32}{MeV} (\cref{fig:17-window}). On its computed highest optical energy, \SI{10.0}{meV}, CsI also leads the polar reach, 3.64 times below GaAs. The same softness gives CsI a poor nuclear edge: from $c_{\mathrm{LA,max}} = \SI{2.46}{km.s^{-1}}$ the edge is \SI{23.1}{keV}, last of the eleven crystals ranked in that row. Two of the four requirements are documented for CsI: crystals far above a gram are grown, and neither caesium-133 nor iodine-127 is radioactive. CsI is slightly hygroscopic, and its operation at millikelvin temperatures after thermal cycling, with a sensor read on it, is not documented. Its entry needs an encapsulated crystal that survives cycling to millikelvin temperatures and a phonon sensor that reads it, deposited on the crystal or on a carrier in thermal contact with it. Its lead then depends on its background: it tolerates $1.56^2 = 2.44$ times TeO$_2$'s background rate density in a background-limited search (\cref{sec:21.7-backgrounds}), and the cosmogenic iodine and caesium isotopes, whose production rates have not been measured for CsI, set that rate.

In the nuclear edge the stiffest crystals lead: diamond leads the survivors on its measured tensor, but in the full pool it shares the leading tier with h-BN, polyethylene and cubic BN, and no leader is resolved. h-BN and polyethylene owe their place to one stiff direction, in the hexagonal plane and along the chains, and their tensors are computed; they reach their edges only for momentum along those directions. LiF leads the polar rate in both pools. The polar reach goes the other way, to soft polar lattices: GaAs among the survivors on its measured longitudinal frequency, CsI in the full pool on a computed one.

The Migdal figure of merit has no leader at independent errors in either pool, for the reason found in \cref{sec:17.6-correlation}, and the light-mediator single-phonon rate ranks no crystals (\cref{tab:17-rank}).

By \cref{sec:20.5}, a leader of an electron-recoil row is read by its charge, through the cascade distribution $P(N_{\mathrm{eh}})$ of \cref{chap:18}, and a leader of a phonon row by its heat, through a collection efficiency that only the detector's kernels fix. A superconducting absorber is read by quasiparticles, with the yield $2\Delta_{\mathrm{BCS}}/E$ below $4\Delta_{\mathrm{BCS}}$ and 0.59 above (\cref{sec:20.4}).

The theory band of a limit (\cref{sec:21.4}) and the ranking of targets answer different questions. An uncertainty that multiplies the rates of every crystal by the same factor, such as the halo or a shared response approximation, moves every limit and no ranking. A ranking moves only through the inputs that differ between crystals, which is why the provenance widths, and not the bands of \cref{tab:21-dominant}, set the probabilities of the guide.

\subsection{Absorption and unranked channels}\label{sec:21.7-other}

Absorption has leaders only mass by mass, the best reach of each decade of \cref{tab:17-abs-decades}. For the dark photon they are zeros of $\varepsilon_1$, where the screening ceiling $L \le 1/(2|\varepsilon_1|)$ is highest: the longitudinal phonon of GaAs, the highest polar branch of sapphire and the plasmon of silicon. All three are survivors. For the axion-like particle with its electron coupling the magnon line of YIG leads from 10~\si{\micro eV} to a meV and clean aluminum from 0.36 to \SI{10.9}{meV}. Neither is a survivor, since the YIG spheres are of millimeter size and no bulk aluminum absorber has been read. The aluminum reach rests on direct pair breaking by the spin current (\cref{sec:15.4-spin}). It holds for a mean free path well above the size of a pair, 1.6 to \SI{2.4}{\micro m}, which only bulk high-purity crystals reach; in a \SI{50}{nm} film the spin channel and the published mapping give comparable reaches.

Magnons, superconductors and the special targets of \cref{chap:16} have no ranking, because no second candidate has the inputs of the same figure of merit. The edge of NiO is set by the magnetic anisotropy that the Heisenberg model leaves out (\cref{sec:14.1}). The scattering rate of aluminum depends more on the gap and mean free path of the absorber actually used than on the choice among loss functions (\cref{sec:15.2}). The one-loop expansion parameter of ZrTe$_5$ is about one, so its rate has no controlled theory at that half-gap (\cref{sec:16.1}). NiO and ZrTe$_5$ lack a gram-scale crystal read at millikelvin temperatures, and aluminum a quasiparticle readout of a bulk absorber.

\subsection{Internal radioactive backgrounds}\label{sec:21.7-backgrounds}

Every figure of merit in the guide ranks signal rates, and a background enters the ranking through a tolerance. Let $b$ be the background rate density in a region of width $\Delta E$. If a search is background limited, the signal must exceed a fixed number of standard deviations of the background count, and its exclusion limit scales as $(b\,\Delta E)^{1/2}/\mathcal{R}$. At the systematics floor, where the background is subtracted only to a fixed fraction, the limit scales as $b\,\Delta E/\mathcal{R}$. Crystal $A$ then keeps its lead over $B$ if $b_A/b_B < (\mathcal{R}_A/\mathcal{R}_B)^2$, or $b_A/b_B < \mathcal{R}_A/\mathcal{R}_B$ at the floor. TeO$_2$ leads CaWO$_4$ in the heavy-mediator rate by 1.22, so it tolerates 1.48 times CaWO$_4$'s background in a background-limited search and 1.22 at the floor. LiF tolerates 3.63 times NaF's background in the polar rate. A lead that rests on a rate factor of order one is reversed by a background ratio of order one. The edges are kinematic and do not move with $b$, but the sensitivity just above them does.

The background that belongs to the crystal itself, and that no shield removes, comes from radioactive nuclei inside it. A $\beta$ emitter spread uniformly through a crystal much larger than the range of its electrons deposits the whole kinetic energy of the electron, and the antineutrino escapes. The spectrum of deposits is then the $\beta$ spectrum, normalized to the activity $\mathcal{A}_{\mathrm{act}}$ per unit mass. For an allowed transition with endpoint $Q_\beta$ it is
\begin{equation}
  n_\beta(T) \propto p_eE_e(Q_\beta - T)^2\,\frac{2\pi y_{\mathrm{S}}}{1 - \ee^{-2\pi y_{\mathrm{S}}}} , \qquad y_{\mathrm{S}} = \frac{\alpha Z_{\mathrm{d}}E_e}{p_e} ,
  \label{eq:21.7-beta}
\end{equation}
with $T$ the kinetic energy of the electron, $E_e = T + m_e$ and $p_e$ its total energy and momentum, $Z_{\mathrm{d}}$ the charge of the daughter nucleus, and the last factor the nonrelativistic Fermi function, which accounts for the Coulomb attraction of the daughter. A unique $n$-th forbidden transition carries in addition the shape factor $\sum_{k=0}^{n}p_e^{2k}q_\nu^{2(n-k)}/[(2k+1)!\,(2n-2k+1)!]$, with $q_\nu = Q_\beta - T$ the momentum of the antineutrino. As $T \to 0$ the factor $p_e$ cancels against the $1/p_e$ of the Fermi function, so the spectrum tends to a constant: the low-energy end of a $\beta$ spectrum is flat. The rate density below \SI{1}{keV} is therefore $\mathcal{A}_{\mathrm{act}}f_\beta$ per keV, with $f_\beta$ the fraction of the spectrum below \SI{1}{keV}. With $Q_\beta = \SI{18.59}{keV}$ for tritium~\cite{iaea2026livechart}, \cref{eq:21.7-beta} gives $f_\beta = 0.0892$. Cosmogenic tritium in silicon is produced at sea level at 124 atoms per kilogram and day. It then contributes 0.169 events per kilogram, day and keV below \SI{1}{keV} after a hundred days and 11.0 at saturation (\cref{app:E.27}). In germanium, at 82 atoms per kilogram and day~\cite{armengaud2017measurement}, the two values are 0.112 and 7.31. Cosmogenic activation is reducible: growth, fabrication and storage underground, and in silicon the out-diffusion of hydrogen, lower it. The other cosmogenic nuclides decay by electron capture, and their L and M lines near 1 and \SI{0.1}{keV} enter the sub-keV region through their L and M fractions. Their production rates have been measured for silicon and germanium~\cite{saldanha2020cosmogenic,armengaud2017measurement} and not for the other survivors.

A lead compound carries a background that no exposure history changes. Its lead contains $^{210}$Pb, half-life 22.2 years, which decays to $^{210}$Bi and, five days later, to $^{210}$Po; in equilibrium the three activities are equal~\cite{iaea2026livechart}. In 84 percent of its decays $^{210}$Pb feeds the level of $^{210}$Bi at \SI{46.5}{keV}, with a $\beta$ endpoint of \SI{17.0}{keV}. The de-excitation of that level, a photon or conversion electrons with their X-rays, is absorbed in the bulk, so these decays deposit at least \SI{46.5}{keV} and never fall below \SI{1}{keV}. The other 16 percent go to the ground state, with $Q_\beta = \SI{63.5}{keV}$. For a daughter as heavy as $^{210}$Bi, $Z_{\mathrm{d}} = 83$, the nonrelativistic factor of \cref{eq:21.7-beta} is crude, and we use instead the traditional relativistic Fermi function of a point charge, with the prefactor 4, evaluated at the nuclear radius $r_A = 1.2A^{1/3}\,\mathrm{fm}$~\cite[Eq.~(6)]{hayen2018high},
\begin{equation}
  F_{\mathrm{Fermi}}(Z_{\mathrm{d}}, T) = 4\,(2p_er_A)^{2(\gamma_{\mathrm{d}} - 1)}\,\ee^{\pi y_{\mathrm{S}}}\,\frac{|\Gamma(\gamma_{\mathrm{d}} + \ii y_{\mathrm{S}})|^2}{\Gamma(1 + 2\gamma_{\mathrm{d}})^2} , \qquad \gamma_{\mathrm{d}} = \sqrt{1 - Z_{\mathrm{d}}^2\alpha^2} ,
  \label{eq:21.7-fermi}
\end{equation}
in place of the last factor of \cref{eq:21.7-beta}, with $y_{\mathrm{S}}$ as there. It gives $f_\beta = 0.0445$. The $^{210}$Bi decay to $^{210}$Po, $Z_{\mathrm{d}} = 84$ and $Q_\beta = \SI{1161}{keV}$, adds $f_\beta = 1.26\times10^{-3}$ with the allowed shape, and the $\alpha$ decay of $^{210}$Po deposits megaelectronvolts. Per becquerel of $^{210}$Pb per kilogram the rate density below \SI{1}{keV} is
\begin{equation}
  \begin{split}
  b_{^{210}\mathrm{Pb}} &= \SI{86400}{s.d^{-1}}\times(0.16 \times 0.0445 + 1.26\times10^{-3})\ \mathrm{keV}^{-1}\times\SI{1}{Bq.kg^{-1}}\\
  &= 724\ \text{per kg, day and keV} .
  \end{split}
  \label{eq:21.7-pb210}
\end{equation}
Of these, 614.7 come from $^{210}$Pb and 109.2, fifteen percent, from $^{210}$Bi. A budget of one event per kilogram, day and keV then needs an activity below \SI{1.38}{mBq} per kilogram of compound: 1.91, 2.23 and \SI{1.60}{mBq} per kilogram of lead in PbSe, PbTe and PbS. The number is an estimate. The relativistic Fermi function raises it by \SI{5}{\percent} over the nonrelativistic 689, mostly through the $^{210}$Bi part, and the first-forbidden transition of $^{210}$Bi, whose shape departs from the allowed one, moves that part further by an amount not computed here. Whether lead of that purity can be obtained and grown into a crystal decides whether the floor of a lead salt is usable, and it is the deciding measurement of the first two rows of \cref{tab:21-guide}.

Primordial nuclei set the only exclusions of the full pool. Natural potassium carries $^{40}$K at \SI{31.7}{Bq} per gram; \SI{89.3}{\percent} of its decays are unique third-forbidden $\beta$ decays with $Q_\beta = \SI{1311}{keV}$~\cite{iaea2026livechart}, for which \cref{eq:21.7-beta} with its shape factor gives $f_\beta = 4.13\times10^{-4}$. KTaO$_3$, \SI{14.6}{\percent} potassium by mass, holds \SI{4.6}{kBq} of $^{40}$K per kilogram and has a $\beta$ rate density of $1.5\times10^5$ per kilogram, day and keV below \SI{1}{keV}. InSb holds \SI{121}{Bq} of $^{115}$In per kilogram, a fourth-forbidden non-unique emitter with $Q_\beta = \SI{497}{keV}$ whose shape is not known; bracketed between the allowed and the unique fourth-forbidden shapes, its rate density lies between $4.8\times10^3$ and $3.9\times10^4$. Both exceed the budget by more than three orders of magnitude, and since only isotopic depletion removes a primordial nucleus, both leave the pool. TeO$_2$ is the opposite case. Its $^{123}$Te is known only to have a half-life above $9.2\times10^{16}$ years~\cite{iaea2026livechart}, so its electron captures are bounded by 693 per kilogram and day, and by the same number per keV if all of them fell in one keV below \SI{1}{keV}. The bound does not decide whether TeO$_2$ meets the budget, and its radioactivity entry stays conditional on a measured capture rate.

Gamma lines of nuclei inside the crystal add a Compton continuum. A photon of energy $E_\gamma$ that scatters once and leaves deposits the electron's recoil energy $T$, and at $T \to 0$ the Klein--Nishina cross-section per electron tends to
\begin{equation}
  \frac{\dd\sigma_{\mathrm{C}}}{\dd T}\bigg|_{T \to 0} = \frac{2\pi r_{\mathrm{cl}}^2m_e}{E_\gamma^2} ,
  \label{eq:21.7-compton}
\end{equation}
with $r_{\mathrm{cl}}$ the classical electron radius; at \SI{1}{keV} the full form for the \SI{1461}{keV} line of $^{40}$K differs from it by $2\times10^{-4}$, and its value there is \SI{1.20e-28}{cm^2.keV^{-1}}. The rate density below \SI{1}{keV} is the photon emission rate times $\nd_e\bar\ell_{\mathrm{esc}}\,\dd\sigma_{\mathrm{C}}/\dd T$ in the thin-target limit, with $\nd_e$ the electron density and $\bar\ell_{\mathrm{esc}}$ the mean path from a point inside the crystal to its surface. For a cube of side $L$ the path from a uniformly distributed interior point in an isotropic direction averages $0.448\,L$, not the mean chord $4V/S = 0.667\,L$ of rays crossing the cube from outside. With $\nd_e = \SI{2e24}{cm^{-3}}$ and $L = \SI{5}{cm}$ a fraction $5.4\times10^{-4}$ of the \SI{1461}{keV} photons deposits in each keV below \SI{1}{keV}. The continuum is flat at low energy like the $\beta$ spectrum, and for an emitter with a $\gamma$ line it adds to the $\beta$ deposits.

\subsection{Open problems}\label{sec:21.7-open}

Three missing inputs limit the guide: the loss functions of the survivors, measured elastic constants, and internal-background rates. The first is the loss function $L(\vect{k},\omega)$ of every survivor near its gap at momenta up to a few keV. Without it the Migdal figure of merit is known only between the plasmon-pole and gap-weighted limits, the prior $X_{\mathrm{pole}}$ stands in for the interband weight, and the row has no resolved leader in either pool.

The second is elastic constants. TeO$_2$, the survivor leader of the heavy-mediator row, and h-BN, polyethylene and cubic BN, which share the leading tier of the full-pool nuclear edge, rest on computed tensors. These carry the method offset that \cref{sec:17.6-correlation} found for the computed family, \SI{4.9}{\percent} on average for silicon, germanium and diamond. For CsI the density of the cold cell is missing.

The third is the internal backgrounds, which the guide has only estimated. The sub-keV rates of $\beta$ emitters use allowed or unique shapes, which are crude for the first-forbidden decay of $^{210}$Bi and unknown for $^{115}$In; the $^{210}$Pb activity of lead that can be grown into a crystal is not known; the $^{123}$Te capture rate is bounded only; and cosmogenic production has been measured for silicon and germanium alone. Until these numbers exist, the radioactivity entries of the guide are conditions, not rates, for every crystal but silicon and germanium.

\subsection{Uncertainty budget of limits and rankings}

\Cref{tab:21-nuisances} is the uncertainty budget that the construction of a limit and of a ranking adds to those of the channels. The halo enters every limit, and near the edge it enters more strongly than any response uncertainty; the statistical construction and the treatment of shared uncertainties move limits by five to forty-five percent; and the rankings rest on width conventions, estimated pair energies and backgrounds that are conditions rather than rates.

{\scriptsize\setlength{\tabcolsep}{3pt}\setlength{\LTcapwidth}{\textwidth}
\begin{longtable}{@{}>{\raggedright\arraybackslash}p{0.25\linewidth}>{\raggedright\arraybackslash}p{0.12\linewidth}>{\raggedright\arraybackslash}p{0.38\linewidth}>{\raggedright\arraybackslash}p{0.09\linewidth}>{\raggedright\arraybackslash}p{0.08\linewidth}@{}}
  \caption{Uncertainty budget of limits and rankings. Classes and bases as in \cref{tab:10-nuisances}. ``Order'' marks an item that acts on a ranking rather than on a rate, ``convention'' a width or a choice fixed by convention. Sizes are for silicon unless stated. The last column gives the section where the item is treated.}
  \label{tab:21-nuisances}\\
  \toprule
  Truncation or input & Type & Size & Basis & Section \\
  \midrule
  \endfirsthead
  \toprule
  Truncation or input & Type & Size & Basis & Section \\
  \midrule
  \endhead
  \bottomrule
  \endfoot
  Reference Earth speed in place of the year average & normali\-zation; tilt in $\mchi$ & $\eta$ within \SI{1}{\percent} to \SI{430.3}{km.s^{-1}}, year average higher by less than \SI{10}{\percent} to \SI{735.6}{km.s^{-1}}; heavy-mediator rate $+\SI{0.9}{\percent}$ at \SI{10}{MeV}, $+\SI{5.6}{\percent}$ at \SI{1}{MeV} & computed & \cref{sec:21.1} \\
  \midrule
  Earth speed at the floor & threshold & June Earth speed lowers the floor by \SI{3.9}{\percent}; a limit uses the exposure average & computed & \cref{sec:21.1} \\
  \midrule
  Halo substructure & normali\-zation; tilt in $\mchi$ & debris flow with 22 percent of the density: $+\SI{8.4}{\percent}$ at \SI{10}{MeV}, $-\SI{16}{\percent}$ at \SI{1}{MeV} & model & \cref{sec:21.1} \\
  \midrule
  Standard halo against every halo & normali\-zation & halo-independent bound weaker by at least 2.00 (\SI{10}{MeV}) and 3.70 (\SI{2}{MeV}), cascade pair bins & computed & \cref{sec:21.2} \\
  \midrule
  Directional response in the halo-independent bound & normali\-zation & program over isotropic formula 1.00--1.13 on the sampled speeds (model anisotropy) & model & \cref{sec:21.2} \\
  \midrule
  Discretized program; final-state assignment & normali\-zation & grid leakage below \SI{0.001}{\percent} at \SI{600}{km.s^{-1}}; misassigned weight at most $0.035\,\varrho^{\mathrm{leak}}_i$ of a bin for a stochastically ordered cascade & computed & \cref{sec:21.2} \\
  \midrule
  Escape tail of the halo & threshold; tilt & edge power 13/2 (standard halo), 15/2 (King), 8--10 (simulated halos); halo fraction within \SI{12}{km.s^{-1}} of $v_{\mathrm{max}}$ from $6.0\times10^{-5}$ to $3.7\times10^{-8}$ & model & \cref{sec:21.3} \\
  \midrule
  Gap or threshold against escape speed at the edge & threshold & \SI{10}{meV} of gap $\equiv -\SI{3.3}{km.s^{-1}}$; limit $\times231$, $\times8.8$, $\times2.7$ at 0.005, 0.010, \SI{0.020}{MeV} above the edge & computed & \cref{sec:21.3} \\
  \midrule
  Year average at the edge & threshold & year-averaged rate over the reference rate 62 at \SI{0.010}{MeV} and 8.5 at \SI{0.020}{MeV} above the reference edge & computed & \cref{sec:21.3} \\
  \midrule
  Theory band of the channel & normali\-zation & silicon, heavy mediator: factor 3.4 (\SI{1}{MeV}), 1.48 (\SI{3}{MeV}), 1.43 (\SI{10}{MeV}); dominant rows per channel in \cref{tab:21-dominant} & computed; model & \cref{sec:21.4} \\
  \midrule
  Statistical construction with a nuisance parameter & normali\-zation & profiled over nominal $\ee^{\sigma_\beta^2N^{\mathrm{up}}}$ with no events, 1.427 at $\sigma_\beta = 0.393$ (Neyman construction, parameter at its conditional best fit); a fixed calibrated threshold charges half the exponent & computed & \cref{sec:21.4} \\
  \midrule
  Excluded cells & normali\-zation & aluminum, $y < 3$ at \SI{1}{keV}: $\times1.25$ (bulk), $\times1.59$ (film); silicon Migdal coherent region: $\times1.49$ at \SI{2}{MeV} & computed & \cref{sec:21.4} \\
  \midrule
  Shared parameter split per exposure & normali\-zation & combined limit $\times1.109$ shared against $\times1.051$ split: about \SI{5}{\percent} too strong & computed & \cref{sec:21.5} \\
  \midrule
  Pair energies outside the validated range & order & logarithmic width 0.21 by convention; lead salts share one tier (0.482, 0.414, 0.102), PbTe 1.000 at $r_{\mathrm{mat}} = 1$; GaSb's estimate beats germanium's measured value with probability 0.83 & convention & \cref{sec:21.7} \\
  \midrule
  Pair energy of germanium at millikelvin & order & \SI{3.37}{eV} from a fit with a fitted phonon-conversion constant, against \SI{2.96}{eV} at \SI{77}{K}; the order against silicon holds for either & est. & \cref{sec:21.7} \\
  \midrule
  $^{210}$Pb in lead compounds & order & 724 per kg, day and keV below \SI{1}{keV} per Bq/kg (relativistic Fermi function, \SI{5}{\percent} above the nonrelativistic value); first-forbidden $^{210}$Bi shape not computed for its \SI{15}{\percent}; activity of available lead not known & est. & \cref{sec:21.7} \\
  \midrule
  Cosmogenic activation of the survivors & order & tritium in silicon 0.169 (100 days at sea level) to 11.0 (saturation) per kg, day and keV below \SI{1}{keV}; production not measured except for silicon and germanium & computed; unsized & \cref{sec:21.7} \\
\end{longtable}}
